\section{Research Objectives and Proposed Work}
\label{sec:proposed_work}

\subsection{Central Hypothesis}
The central hypothesis of this project is that integrating genomic, epigenomic, metabolomic, and peripheral immunomic data across publicly available LBD cohorts will reveal convergent multi-layer molecular signatures that: (a) discriminate LBD from AD and healthy aging with higher accuracy than any single omics layer alone; and (b) identify biologically coherent, druggable pathway perturbations that underlie LBD pathogenesis.

\subsection{Specific Aims}

\subsubsection*{Aim 1 -- Multi-Omics Biomarker Discovery}
Identify a minimal multi-omics biomarker signature in blood-accessible samples (plasma metabolomics, peripheral immune cell transcriptomics, and blood-derived methylation data) that can discriminate (i) DLB from AD, (ii) LBD subtypes (DLB vs. PDD) from each other, and (iii) early from late disease stages. This aim directly addresses the diagnostic gap identified in the literature review and builds on the metabolomics \cite{frontiers2024metabolomics}, single-cell immune \cite{singlecell2024immune}, and biomarker review \cite{boiten2025biomarkers} findings.

\subsubsection*{Aim 2 -- Cross-Layer Pathway Convergence Analysis}
Identify biological pathways that show coordinated dysregulation across two or more omics layers. Priority pathways---lysosomal biology, synaptic vesicle trafficking, sphingolipid metabolism, and neuroinflammation---are nominated by prior single-layer studies \cite{chia2025gwas, nabais2024methylation, frontiers2024metabolomics}. The aim is to formally test pathway convergence using network-based multi-omics integration and to rank pathways by the number of independently dysregulated omics features they contain.

\subsubsection*{Aim 3 -- Drug Target Prioritisation}
Apply Mendelian randomisation and network centrality analysis to prioritise druggable targets within convergent pathways. Candidate targets will be cross-referenced against the Open Targets Platform to assess existing drug-target evidence, clinical trial status, and safety profiles. This aim positions the project output for direct experimental follow-up or clinical trial hypothesis generation.

\subsection{Innovation}
The proposed project is innovative in three respects. \textbf{First}, while each omics layer has been studied in LBD individually, no study to date has computationally integrated all five layers (genomics, epigenomics, transcriptomics, metabolomics, immunomics) in a single analytical framework for LBD specifically. \textbf{Second}, the emphasis on blood-accessible omics layers ensures that findings have direct translational potential as minimally invasive clinical tests. \textbf{Third}, the project is designed to be fully reproducible: all data sources are open-access, all analysis code will be version-controlled and publicly released, and the pipeline will be containerised (Docker/Singularity) to enable external replication.

\subsection{Scope and Scale}
This is conceived as a \textbf{computational, data-driven project}. No new sample collection, wet-laboratory experiments, or patient recruitment are required for the primary aims. The project is therefore well-suited for a single researcher or a small team with access to standard computational infrastructure (a workstation or cloud computing account). The expected duration for a thorough completion is 12--18 months for a researcher with intermediate bioinformatics skills, or 6--9 months for an experienced computational biologist.
